\chapter{Cross-Layer Lessons and Future Outlook}
\label{ch:cross-layer}

\section{Summary of Contributions}

This dissertation set out to bridge the gap between the theoretical promises of quantum-resilient cryptography and the demanding realities of datacenter deployment. By tracing the path of secure communication from the physical generation of entangled photons up to the observability of distributed AI workloads, it has demonstrated that scalable security cannot be achieved within a single layer of the network stack.

The contributions directly address the four research questions posed in Chapter~\ref{ch:introduction}:
\begin{itemize}
    \item \textbf{RQ3 (Integrated Photonic Sources):} Chapter~\ref{ch:quantum-sources} established that DCM resonators in silicon nitride can serve as stable, predictable sources of frequency-bin entangled photons. By developing a digital twin and experimentally validating the thermal stability of the FSR, the work showed how complex device physics can be abstracted into reliable network components.
    \item \textbf{RQ2 (QKD Architecture):} Chapter~\ref{ch:qkd-dpu} addressed the computational bottleneck of QKD post-processing. By offloading the entire classical pipeline---including LDPC error correction and privacy amplification---onto the Arm cores of a BlueField-3 DPU, the architecture isolates key generation from the host CPU and enables line-rate quantum-secured RDMA.
    \item \textbf{RQ1 (PQC Deployment):} Chapters~\ref{ch:dpu-pqc} and~\ref{ch:ai-pipelines} demonstrated how post-quantum cryptography can be efficiently deployed on programmable network infrastructure. By partitioning heavy cryptographic workloads across CPU, GPU, and DPU substrates, the research showed that PQC offload can preserve host and accelerator capacity while still meeting the deadlines of latency-sensitive, end-to-end AI pipelines.
    \item \textbf{RQ4 (Cross-Layer Interactions):} Chapters~\ref{ch:rdma} and~\ref{ch:confidential-ai} revealed how choices at one layer constrain or enable others. High-throughput optical fabrics expose timing, ordering, congestion, and routing behaviours that can undermine cryptographic guarantees, while cross-layer observability can support privacy-preserving detection once confidential AI payloads become opaque.
\end{itemize}

Taken together, these contributions show that practical quantum resilience requires co-designing the device, protocol, and infrastructure layers.

\section{Cross-Layer Lessons}

The three layers developed across this thesis are not independent. Their dependencies
form reciprocal constraints that, taken together, answer RQ4 and shape system
design:

\begin{itemize}
    \item \textbf{Device to Protocol:} Achieving a high SKR for any QKD protocol begins at the device layer, by ensuring spectral purity and FSR stability. For instance, the experimental validation in Chapter~\ref{ch:quantum-sources} demonstrated that the dual-coupled microring source maintains FSR stability under thermal tuning, which is critical for aligning with standard telecom grids. If this stability drifts, the resulting basis mismatches directly inflate the QBER. As shown in Chapter~\ref{ch:qkd-dpu}, this QBER dictates the computational burden of the protocol layer: higher error rates require larger LDPC matrices, lower code rates, and more belief-propagation iterations, increasing the processing latency and memory footprint during QKD post-processing.
    
    \item \textbf{Protocol to Infrastructure:} The computational complexity of the cryptographic protocol dictates the required infrastructure. Chapter~\ref{ch:dpu-pqc} showed that code-based and lattice-based PQC workloads place different pressure on host CPUs, DPU Arm cores, and embedded GPU acceleration: memory-bound HQC benefits from isolated DPU execution under contention, while batched ML-KEM and ML-DSA benefit from the parallelism and shared memory domain of an embedded GPU. Chapter~\ref{ch:ai-pipelines} then showed that this placement question survives at application scale, where terminating the secure receive path on the DPU preserves GPU capacity and keeps a real-time image pipeline within its frame budget. The protocol's mathematical structure thus forces a shift in datacenter architecture, moving cryptographic execution out of the host and into the network fabric where latency, batching, isolation, and accelerator contention can be controlled.
    
    \item \textbf{Infrastructure to Protocol:} Conversely, the physical and operational realities of the infrastructure can undermine the mathematical guarantees of the protocol. Chapter~\ref{ch:rdma} demonstrated that the ordering, timing, congestion-control, and routing mechanisms of RDMA over high-throughput optical fabrics expose fabric-level structural channels that a network-level adversary can exploit, even when the payloads are cryptographically secured. Defending against these infrastructure-derived threats requires infrastructure-level visibility. As explored in Chapter~\ref{ch:confidential-ai}, once payloads become opaque due to robust encryption, defenders must rely on cross-layer observability---correlating network telemetry with hardware performance counters---to detect structural inconsistencies in confidential AI training clusters.
\end{itemize}

This cross-layer perspective is the reason the dissertation combines topics---photonics,
cryptography, and datacenter networking---that would otherwise appear separate. The unifying
question that runs through every chapter is how secure communication systems can be made
scalable and deployable when the binding constraints are distributed across physical,
mathematical, and operational boundaries.

\section{Limitations}

While this dissertation spans multiple layers of the network stack, its findings are subject to several practical and experimental boundaries:

\begin{itemize}
    \item \textbf{Experimental Scope:} Although the thesis bridges theoretical protocols and deployable infrastructure, many of the evaluations were conducted in controlled testbeds rather than live, multi-tenant production datacenters. For example, the QKD post-processing pipeline in Chapter~\ref{ch:qkd-dpu} was evaluated using a synthetic BB84 source that injected controlled QBER profiles. While this isolated the computational performance of the DPU from the environmental drift, dark counts, and afterpulsing that plague physical optical hardware, it abstracts away the messy realities of a live quantum channel.
    \item \textbf{Hardware Specificity:} The architectural solutions proposed are highly optimised for specific hardware platforms. The infrastructure-layer offloading relies heavily on the capabilities of the NVIDIA BlueField-3 DPU, particularly its Arm cores, embedded acceleration options, and inline cryptographic engines. Similarly, the device-layer findings in Chapter~\ref{ch:quantum-sources} are specific to silicon nitride ($\text{Si}_3\text{N}_4$) photonic platforms. The underlying principles generalise, but the exact performance metrics and integration strategies are tightly coupled to these technologies.
    \item \textbf{Technology Maturity:} The landscape of quantum-resilient cryptography is evolving rapidly. The PQC algorithms evaluated in this thesis represent the current state of NIST standardisation, but the field remains active, and future algorithmic tweaks or entirely new schemes may alter the computational bottlenecks identified here. Likewise, integrated quantum sources are still maturing from laboratory demonstrations to commercial, telecom-grade components.
\end{itemize}

\section{Future Outlook}

The transition to quantum-resilient communication is not a single upgrade but an ongoing architectural evolution. Building on the findings of this thesis, several critical directions emerge for future research and deployment:

\begin{itemize}
    \item \textbf{End-to-End QKD in WDM (TU/e Lab):} The integrated quantum source characterised in Chapter~\ref{ch:quantum-sources} is currently serving as the foundation for ongoing experimental work at the Eindhoven University of Technology (TU/e) QKD lab. The objective of this active research is to construct a complete, end-to-end QKD system that operates concurrently with classical traffic over WDM. Achieving stable coexistence between fragile single-photon quantum channels and high-power classical data streams on the same fibre infrastructure remains one of the most significant hurdles for wide-scale QKD deployment.
    
    \item \textbf{AI-Driven Datacenter Evolution:} Datacenters are advancing at an unprecedented rate, driven by the explosive growth of AI workloads, including large-scale inference and agentic AI systems. This rapid evolution of the infrastructure layer demands security architectures that can scale and adapt just as fast. The static security perimeters of the past are insufficient for the dynamic, high-bandwidth, and latency-sensitive demands of modern AI clusters.
    
    \item \textbf{Cryptographic Agility and Vigilance:} The cryptographic community must remain vigilant and agile. While lattice- and hash-based PQC algorithms currently offer the best defence against quantum threats, they are not mathematically proven to be unconditionally secure. A cryptanalytic breakthrough against their underlying hardness assumptions could break current standards at any time. Active research into alternative cryptographic primitives is essential, and network infrastructure must be designed with absolute cryptographic agility, allowing systems to seamlessly swap algorithms without requiring hardware replacements.
    
    \item \textbf{Multi-Hybrid Architectures:} To hedge against the risk of algorithmic failures, future work should focus on multi-hybrid key exchange architectures. By integrating the information-theoretic security of QKD with multiple software deployable PQC algorithms on programmable network hardware (such as DPUs), networks can achieve layered defence. If one cryptographic mechanism falls, the others maintain the security of the channel.
    
    \item \textbf{Confidential AI Observability:} As AI training clusters become increasingly complex and multi-tenant, the attack surface expands. Future research must advance cross-layer intrusion detection and observability. When payloads are fully encrypted and hardware is shared among competing tenants, defending these clusters will require sophisticated, AI-driven telemetry analysis that can detect structural anomalies and fabric-level leakage in real time, preserving the confidentiality of the models and data without degrading training performance.
\end{itemize}

The quantum threat has forced a fundamental re-evaluation of how digital trust is established and maintained. By embracing a cross-layer approach that spans from the physics of photonics to the architecture of the datacenter, the next generation of secure communication systems can be built to withstand both the quantum computers of tomorrow and the scaling demands of today.
